% 分野別類題: ケーリー・ハミルトンの定理と行列多項式
% 固有多項式による次数下げ・多項式の割り算を用いた A^n の計算 3 問（2x2 基本2問＋3x3 応用1問）。
% コンパイル: TEXINPUTS="../../../00_common:$TEXINPUTS" latexmk -xelatex problems.tex
\documentclass[a4paper,11pt]{article}
\usepackage{preamble}

\begin{document}

\kamokumei{類題演習 --- ケーリー・ハミルトンの定理と行列多項式}

\shijibun{以下の各問に答えよ。ケーリー・ハミルトンの定理（$A$ の固有多項式 $\varphi(\lambda) = \det(\lambda I - A)$ に対し $\varphi(A) = O$ が成り立つ）を用いて次数下げを行い、計算過程を示すこと。}

\shomon{問1.}{行列 $A = \begin{pmatrix} 1 & 2 \\ 2 & 1 \end{pmatrix}$ について、ケーリー・ハミルトンの定理を用いて $A^{5} - A^{2} - 16A$ を計算せよ。}

\shomon{問2.}{行列 $B = \begin{pmatrix} 2 & 1 \\ 0 & 3 \end{pmatrix}$ について、$\lambda^{6}$ を $B$ の固有多項式で割った余りを用いて $B^{6}$ を求めよ。}

\shomon{問3.}{行列 $C = \begin{pmatrix} 2 & 0 & 0 \\ 0 & 1 & 1 \\ 0 & -2 & 4 \end{pmatrix}$ について、ケーリー・ハミルトンの定理を用いて $C^{4}$ を計算せよ。}

\end{document}
